\section{Removing one integer witness from the joint-sign family}
\label{int-section}

The smallest-degree branch and the smallest-variable branch need different
comparisons.  Here we improve four recent equations over \emph{unrestricted
integers}.  The principal endpoint is a single polynomial in eighteen
integer variables, of exact degree $88$, with $105$ live binary arithmetic
operations.  Its direct predecessor in the incoming joint-guards packet
uses nineteen variables, degree $88$, and $109$ operations.

\begin{center}
\begin{tabular}{rrrrrr}
\toprule
\multicolumn{3}{c}{Incoming predecessor}&\multicolumn{3}{c}{This composition}\\
Variables&Degree&Operations&Variables&Degree&Operations\\
\midrule
19&88&109&18&88&105\\
20&32&112&19&32&108\\
21&20&115&20&20&111\\
22&16&118&21&16&114\\
\bottomrule
\end{tabular}
\end{center}
Every entry in this table has domain $\Z$; none is obtained by silently
relabelling positive witnesses.  All four new equations retain precisely
the same six fixed numerals as their direct predecessors.  These are new
compositions and measured endpoints, not a new general elimination lemma:
source~73, \emph{Guard reanchoring}, already uses the complement coordinate
and monic transport elimination proved below.  Its application there is to
an earlier forty-one-variable system.  The incoming joint-sign projection
family retains the eliminable transport coordinate, which lets the same
method produce the four improvements in the table.

\subsection{An exact integer-polynomial projection}
\label{int-projection}

Suppose a source system contains polynomial abbreviations $q,W$ independent
of two integer coordinates $F,\alpha$, and put $C=q-F-\alpha$.
Suppose one of its rows is
\begin{equation}
 (K_0+W)C-F-z(q-1)=0.                         \label{int-old-transport}
\end{equation}
First replace $\alpha$ by the free integer coordinate $C$.  This is an
invertible polynomial change of coordinates, with inverse
$\alpha=q-F-C$.  The row is now monic in $F$.  Define
\begin{equation}
 \widehat F=(K_0+W)C-z(q-1),\qquad
 \widehat\alpha=q-\widehat F-C.               \label{int-restore}
\end{equation}
Substitute these polynomials in every other row and delete
Equation~\eqref{int-old-transport}.

\begin{proposition}
\label{int-projection-thm}
The old and new integer zero sets are in bijection.  The projection fixes
every other coordinate and sends $(F,\alpha)$ to $C=q-F-\alpha$;
its inverse is Equation~\eqref{int-restore}.
\end{proposition}
\begin{proof}
At an old zero the transport row gives $F=\widehat F$, and then
$\alpha=\widehat\alpha$.  All substituted rows vanish.  Conversely,
restore these two values at a new zero.  The deleted row vanishes
identically, while every retained row vanishes by the defining
substitution.  Both compositions fix the original coordinates, because
the old transport equation determines $F$ uniquely.  There is no division,
positivity assumption, coprimality assumption, or new choice of canonical
Pell witnesses in this argument.
\end{proof}

The rank expression also becomes shorter:
\begin{align}
 &\bigl((q-1)(q-F)+\alpha\bigr)(q^2-1)
       +(M_C+qM_F)J \notag\\
 &\hspace{15mm}=
   \bigl(q(q-\widehat F)-C\bigr)(q^2-1)
       +(M_C+qM_F)J.                         \label{int-rank-identity}
\end{align}
This is a literal polynomial identity after the substitution, not merely
an equality at solutions.  The determinant guard remains valid with
$F=\widehat F$ because it is substituted everywhere, including its sign
constraints.  No coefficient port is replaced or borrowed from another
construction.

\subsection{The complete eighteen-variable polynomial}
\label{int-polynomial}

The eighteen unknowns are exactly
\begin{equation}
 J_3,C,z,f,h,r,v,w,\tau,y,\gamma,o,j,c,t,p_1,p_2,p_3\in\Z.
                                                        \label{int-vars}
\end{equation}
The variable $C$ is unrestricted; its required sign is enforced by the
polynomial's inherited guard.  All following symbols are polynomial
abbreviations and do not add quantified variables.  Define
\begin{align*}
 J&=3J_3+1,&q&=\beta J+1,\\
 W&=4(3w+1),&X&=qW,\\
 Y&=(2v+1)q^3,&F&=(K_0+W)C-z(q-1),\\
 E&=XY,&a&=E+Y,\\
 H&=4a+3,&\Delta&=a^2+H,\\
 V&=fo-c,&I&=(2r+1)c^2,\\
 k&=1-V+jc+hE,&\eta&=c-kY,\\
 \zeta&=k-\eta,&T&=\eta\zeta.
\end{align*}
Then put
\begin{align*}
 R&=\bigl(q(q-F)-C\bigr)(q^2-1)+(M_C+qM_F)J,\\
 d&=(A F)(B C)-(q-AF-BC)^2,\\
 P_+&=1+p_1^2+p_2^2+p_3(p_3+1).
\end{align*}
Here $P_+\ge1$ for every integer triple; its standard surjectivity onto
positive integers is used only in the inherited joint-sign theorem.
The six residuals are
\begin{align}
 L_1&=\tau^2-EkY(k+EkY)-1,\notag\\
 L_2&=(X+ac+\gamma H)^2-\Delta c^2-1,\notag\\
 L_3&=cj-V-R,\notag\\
 L_4&=I^2-\Delta(f^2-1),\notag\\
 L_5&=I^2(V^2-y^2)+y^2-1,\notag\\
 L_6&=dT-(t^2-d-T+1)^2-P_+.                    \label{int-six-rows}
\end{align}
The complete equation is
\begin{equation}
 \boxed{\quad Q_{18}=L_1^2+L_2^2+L_3^2+L_4^2+L_5^2+L_6^2=0,
          \qquad\mathbf x\in\Z^{18}.\quad}    \label{int-Q18}
\end{equation}
The JSON uses the names \texttt{v\_scale}, \texttt{w3},
\texttt{pair\_root}, and \texttt{pair\_pos1--3} for $v,w,t,p_1,p_2,p_3$.
Those are only spelling changes.  Its $105$-gate straight-line program is
an explicit division-free definition of the same ordinary polynomial.

\subsection{The six fixed numerals}
\label{int-coefficients}

The tuple is the complete source-34 native binding named
\texttt{mixed-parity-unrounded-1064921}, together with the two determinant
scales.  In the displayed equation,
\[
 \beta=2^{1064921}-1,\qquad A=2^{1113},\qquad B=2^{666104}.
\]
The other three fixed integers $K_0,M_C,M_F$ are supplied as literal,
compressed, minimal unsigned little-endian byte strings, with their
complete finite construction retained alongside them.  In particular
$M_F$ is the \emph{already shifted} source port.  It is not shifted again.
All six numeral payloads together occupy $615883$ bytes before compression.

For clarity, the finite construction of the native ports is as follows.
The retained literal $215$-state table determines $943$ ordered Wang tiles
by the explicit finite map in
\path{integer/upstream/guards/upstream/source34/source34_builder.py}.
With radix $\rho=2^{121}$ and block length $8801$, that map builds a
$989$-element support $\mathcal E$, its explicit nonnegative coefficients
$c_e$, and the dictionary
\[
 d_s=\mathbf1_{s\in\{40,41,1100,2200\}}
      +\sum_{e\in\mathcal E}c_e
          \bigl(\mathbf1_{s=4393-e}+\mathbf1_{s=5495-e}\bigr).
\]
Here the indicator sum counts coincident contributions.  Put
$D_C=\sum_s d_s\rho^s$ and $\mathcal B=\rho^{8801}$.
The finite recipes then give
\begin{align*}
 K_0&=D_C+\mathcal B\rho,\\
 M_C&=\mathcal B-1
         -\sum_{e\in\mathcal E\setminus\{1\}}\rho^e
         -2\rho^{1063},\\
 M_F&=(\mathcal B-1)+4+
      \mu(\rho^{4393}+\rho^{5495})
      +\sum_{j=0}^{9}(\rho^{1000+2j}+\rho^{1021+2j})\\
     &\hspace{8mm}+\rho^{2195}+\rho^{3296}
       +(2^{119}-1)\sum_{e=8792}^{8798}\rho^e
       +(2^{27}-1)\rho^{8799},\\
 \mu&=1022+\sum_{j=0}^{43}2^{10+2j}.
\end{align*}
The finite support, every tile and coefficient, and the generating table
are included as data and executable finite recipes.  This is not an
appeal to an unspecified universal polynomial or a free compiler index.
The verifier reconstructs the entire tuple, compares its integer values
to the literal bytes, and replays $69230$ finite convolution-basis checks.
These finite checks bind the numerals; the unbounded semantic theorem is
a separate inherited argument.

\subsection{Why the joint row really enforces two positive signs}

The inherited guard has a useful elementary proof, which also explains why
the eighteen displayed inputs may range over all integers.  For integers
$x,y,t$ put
\[
 G(x,y,t)=xy-(t^2-x-y+1)^2.
\]
Then $x,y>0$ if and only if $G(x,y,t)>0$ for some integer $t$.
Indeed, $G>0$ implies $xy>0$.  If $x=-a,y=-b$ with $a,b>0$, then
$(t^2+a+b+1)^2>ab$, a contradiction.  Conversely, for positive $x,y$ put
$n=x+y-1$.  We have $xy\ge n$.  If $n$ is a square, take $t^2=n$.
Otherwise write $m^2<n<(m+1)^2$ and choose whichever of $m,m+1$ makes
$|t^2-n|$ smaller.  This distance is at most $m$, so
$(t^2-n)^2\le m^2<n\le xy$.

Every positive integer is a value of
$1+p_1^2+p_2^2+p_3(p_3+1)$.  This follows from the three-square theorem
applied to $4n-3$: exactly one square is odd modulo four, giving the
displayed expression after division by four.  Consequently
\[
 x,y>0\quad\Longleftrightarrow\quad
 \exists t,p_1,p_2,p_3\in\Z\quad
 G(x,y,t)=1+p_1^2+p_2^2+p_3(p_3+1).
\]
Apply this to $x=d,y=T$ in $L_6$.  The new projection substitutes $F$
in both occurrences of the determinant and in the rank; it does not
delete this guard or infer positivity from the spelling of a variable.
The further interpretation of $d>0$ and $T>0$ is exactly the retained
parent's determinant and signed-Pell argument.

\subsection{Correctness and its precise proof boundary}
\label{int-correctness}

Let $Q_{19}^{\mathrm{old}}$ be the exact incoming polynomial
\path{joint/joint19_degree88.json} in the archive
\path{ZFC_Diophantine_Joint_Guards_2026-10-24.zip}.
The repository is pinned at commit \texttt{c1a102936aca};
the complete identifier appears on the title page.
Its source graph and the exact parent-to-polynomial binding are retained.
The preceding projection gives the unconditional identity
\[
 Q_{18}(\mathbf x)=
 Q_{19}^{\mathrm{old}}(\mathbf x,
       F=\widehat F,\alpha=\widehat\alpha),
\]
where the input correspondence replaces the old $F,\alpha$ by the new
$C$.  The deleted squared residual is identically zero.
Consequently
\begin{equation}
 \exists\mathbf x\in\Z^{18}\ Q_{18}=0
 \quad\Longleftrightarrow\quad
 \exists\mathbf y\in\Z^{19}\ Q_{19}^{\mathrm{old}}=0.       \label{int-bijection}
\end{equation}
This equivalence is proved over arbitrary integer values of the six
coefficient ports, so there is no hidden numerical or sign assumption
in the new step.

For the authentic tuple the retained parent proof proceeds through the
joint-positivity lemma, determinant sign restoration, the signed-power
bootstrap, native Wang-tableau compilation, Pell calibration, and the
literal machine/refutation theorem.  Composing that ordinary proof chain
with Equation~\eqref{int-bijection} yields the intended theorem
\[
 \ZFC\ \vdash\
 \bigl(\neg\exists\mathbf x\in\Z^{18}\ Q_{18}(\mathbf x)=0\bigr)
       \ \longleftrightarrow\ \Con(\ZFC).
\]
This statement neither assumes nor proves consistency.  The new
projection is an elementary argument formalizable in ZFC.  The archive
does not provide a new end-to-end Lean/Rocq certificate or an exported
object-language ZFC derivation: it inherits the named unbounded
calibration interfaces, their ordinary proofs, and their existing formal
verification limitations.  A finite replay cannot replace those interfaces.

\subsection{Exact degree and operation count}
\label{int-degree}

Give each quantified coordinate degree one and the fixed ports degree
zero.  The six residual degrees in Equation~\eqref{int-six-rows} are
\[
 (34,14,5,14,10,44).
\]
In particular, the rank row rises from degree four to degree five after
the projection; this does not change the dominant degree.  Write
$X_*,Y_*,E_*,\eta_*$ for the top homogeneous parts.  They are
\begin{align*}
 X_*&=36\beta J_3w,&Y_*&=54\beta^3J_3^3v,\\
 E_*&=X_*Y_*,&\eta_*&=-hX_*Y_*^2.
\end{align*}
Thus $T$ has degree $22$ and top part $-\eta_*^2$, whereas $d$ has
degree at most four.  Only $L_6^2$ reaches degree $88$, with top part
$\eta_*^8$.  In particular
\begin{equation}
 [h^8w^8v^{16}J_3^{56}]Q_{18}
       =2^{32}3^{64}\beta^{56}>0.                    \label{int-top}
\end{equation}
This proves exact degree on the actual fixed-coefficient slice, rather
than merely giving a degree propagated through a circuit.

The delivered DAG uses $54$ multiplications, $27$ additions, and $24$
subtractions.  Every square and every scalar multiplication is charged;
fixed numeral leaves and reuse of a register are free.  The six residuals
need $94$ operations, and six squarings plus five additions give $105$.
The predecessor needs $109$.  Restoring $F$ costs additional arithmetic,
but eliminating its transport row and shortening the rank expression
save four live operations overall.  No minimum-operation claim is made.

\subsection{Three companion tradeoffs}
\label{int-companions}

The same transformation applies to the incoming $20/32$, $21/20$, and
$22/16$ joint-sign polynomials.  Their three resulting packets have,
respectively, nineteen, twenty, and twenty-one unrestricted coordinates.
They retain some intermediate coordinates that Equation~\eqref{int-Q18}
eliminates.  Precisely:
\begin{itemize}
\item The nineteen-variable, degree-$32$ form retains $\eta$ and the row
      $c-\eta-kY=0$.
\item The twenty-variable, degree-$20$ form also retains $E$ and the row
      $E-XY=0$, while $a=E+Y$ remains an abbreviation.
\item The twenty-one-variable, degree-$16$ form also retains $a$ and
      $a-E-Y=0$.  Its source rows use $a-E$ in place of $Y$ where specified
      by the parent.  It retains the parent's reduced auxiliary residual
      $\Delta(f^2-1)(V^2-y^2)+y^2-1$.
\end{itemize}
The last replacement is justified in the parent by retaining
$I^2-\Delta(f^2-1)$ and subtracting its multiple $V^2-y^2$ from the
original auxiliary row.  It is a zero-set-preserving row operation, not
an unrestricted polynomial identity.  Our new substitution preserves the
particular parent polynomial for each endpoint.

Their independently checked degree certificates are
\begin{align*}
 [\eta^4h^4w^4v^4J_3^{16}]Q_{19}&=2^{12}3^{20}\beta^{16},\\
 [f^4v^4J_3^{12}]Q_{20}&=2^43^{12}\beta^{12},\\
 [a^4f^8o^4]Q_{21}&=1.
\end{align*}
All are positive for the authentic tuple.  The expanded row degrees and
literal circuit counts are in the verification receipts.  These exact
integer-domain improvements coexist with lower-variable, much-higher-degree
equations; no scalar measure makes every tradeoff comparable.

\subsection{Reproduction and provenance}
\label{int-reproduction}

The component commands are \texttt{python3 build.py},
\texttt{python3 verify.py}, and \texttt{python3 -O verify.py}.
The independent checker does not import the new builder.  It authenticates
the retained source files and coefficient tuple, rechecks the parent DAGs
against their literal polynomials, proves every new residual identity by
exact sparse algebra over validated cut registers, checks degree bounds
and nonzero top coefficients, verifies the sum-of-squares finalizer,
and rejects an operation-preserving sign corruption.  Normal and optimized
Python runs pass.  No diagnostic search for a consistency-equation zero
is used.

The primary polynomial is \path{transport18_degree88.json}; the three
companions are \path{transport19_degree32.json},
\path{transport20_degree20.json}, and
\path{transport21_degree16.json}.  Each declares every witness, fixed
numeral, gate, residual, and restoration map.  The inherited proof chain,
source~73's earlier method, literal finite recipes, fixed bytes, licenses,
and source hashes are retained under \texttt{upstream/}.  The method is
credited to that earlier report; the current contribution is its exact
composition with the newer joint-sign family.
